% Part: proof-theory
% Chapter: natural-deduction
% Section: sequents

\documentclass[../../../include/open-logic-section]{subfiles}

\begin{document}

\olfileid[tr]{pt}{nat}{seq}

\olsection{Sekantlarla Doğal Tümdengelim: \Log{N2c} ve~\Log{N2i}}

\Log{N1c} ya da \Log{N1i}'deki bir !!^a{proof}~$\delta$, son-!!{formula}{}ü
$!A$'nın açık varsayımlarından çıktığını gösterir. $\Gamma$, $!B^x$'in
$\delta$ içinde açık bir varsayım olarak geçtiği !!{formula}s~$!B$ kümesiyse
bunu $\delta \Proves \Gamma \Sequent !A$ biçiminde yazabiliriz. Bu, sekantlar
için de bir doğal tümdengelim sistemi formüle edebileceğimizi düşündürür.
Böyle bir sistemde her sekant, sağdaki !!{formula}{}ün hangi açık varsayımlara
dayandığı, yani bir sekantla biten alt-!!{proof} içinde hangilerinin açık
olduğu bilgisini taşır. Bu, sekantlarla doğal tümdengelim ya da ``sekant
tarzında'' doğal tümdengelimdir; ortaya çıkan sistemlere \Log{N2c}
ve~\Log{N2i} diyoruz.

\Log{N1} ile \Log{N2} arasındaki karşılıklılığı daha yakın kılmak için soldaki
$\Gamma$ kümelerini etiketli !!{formula}s~$x:!B$'den oluşturalım. \Log{N1}'deki
kurallar yalnızca öncüllerdeki !!{formula}s, yani son-!!{formula}s üzerinde
işlem yapar; söz konusu son formüller doğrudan alt-!!{proof}s{}a aittir.
Dolayısıyla \Log{N2}'nin kuralları da
sekantın yalnızca ardılı üzerinde işlem yapar. Böylece sekantların solundaki
etiketli formüllere kısaca bağlam diyebiliriz.

\Log{N2} ispatlarında varsayımlar yerine en üstte şu biçimde aksiyom
sekantları bulunur:
\[x:!A \Sequent !A.\]
Kurallar \Log{N1}'in kurallarına tam olarak karşılık gelir ve
\olref{tab:N2}'de verilmiştir.

\subfile{rules-N2}

(\Log{N1c} ve \Log{N1i}'de yaptığımız gibi) \Intro{\lif}, \Intro{\lnot},
\Elim{\lor} ve \Elim{\lexists} kurallarının varsayımları boşaltmasına izin
vermekle birlikte bunu zorunlu tutmadığımızdan, karşılık gelen bağlam
!!{formula}s{}i $!A$, $!B$ ve $!A(a)$ ilgili öncülün bağlamında bulunmadığında
da \Log{N2c} ve~\Log{N2i}'deki karşılık gelen kuralların doğru uygulanmasına
izin veririz.

\begin{prop}
$\delta$, \Log{N1c} ya da \Log{N1i}'de $!A$'nın bir !!a{proof}{}ı ise
\Log{N2c}'de (\Log{N2i}'de) $\Gamma \Sequent !A$'nın bir $\delta'$ !!a{proof}{}ı
vardır; burada $\Gamma$, $!B^x$'in $\delta$'nın açık bir varsayımı olduğu tüm
$x:!B$'lerin kümesidir.
\end{prop}

\begin{proof}
  $n=\pheight{\delta}$ üzerinde tümevarımla. $n=0$ ise $\delta$ tek bir açık
  $!A^x$ varsayımıdır. Bu durumda $\Gamma = \{x:!A\}$ ve $\Gamma \Sequent !A$,
  \Log{N2c} ya da~\Log{N2i}'nin bir aksiyomudur.

  $n>0$ ise $\delta$'daki son çıkarıma göre durumları ayırırız. Yalnızca bir
  durumu ele alıyoruz; geri kalanı alıştırma olarak bırakılmıştır.

  Son çıkarımın $\Intro{\lif}$ olduğunu varsayalım. Bu durumda $\delta$ şöyle
  biter:
  \[
  \AxiomC{$\Discharge{!B}{x}$}
  \RightLabel{$\delta_1$}
  \DeduceC{$!C$}
  \DischargeRule{\Intro{\lif}}{x}
  \UnaryInfC{$!B \lif !C$}
  \DisplayProof
  \]
  Tümevarım varsayımı gereği \Log{N2c}'de (\Log{N2i}'de) $\Gamma' \Sequent
  !C$'nin bir $\delta_1'$ !!a{proof}{}ı vardır; burada $\Gamma'$, $\delta_1$
  içinde açık olan etiketli varsayımlardır. $!B^x$, $\delta_1$'in açık bir
  varsayımıysa \Intro{\lif} çıkarımında boşaltılır ve $\delta$'nın açık
  varsayımları $\Gamma = \Gamma' \setminus \{x:!B\}$ olur. Başka bir deyişle
  $\delta_1'$, $x:!B, \Gamma \Sequent !C$'nin bir !!a{proof}{}ıdır ve
  $\delta'$ olarak şunu alabiliriz:
  \[
  \AxiomC{}
  \RightLabel{$\delta_1'$}
  \Deduce$x:!B, \Gamma \fCenter !C$
  \RightLabel{\Intro{\lif}}
  \UnaryInf$\Gamma \fCenter !B \lif !C$
  \DisplayProof
  \]
  $!B^x$, $\delta_1$'in açık bir varsayımı değilse \Intro{\lif} bir varsayımı
  boşaltmaz ve $\delta$ ile $\delta_1$'in açık varsayımları aynıdır.
  \Log{N2c} ve~\Log{N2i}'deki \Intro{\lif} öncülünde $x:!B$ bağlam
  formülünün bulunması gerekmediğinden $\delta'$ olarak şunu alabiliriz:
  \[
  \AxiomC{}
  \RightLabel{$\delta_1'$}
  \Deduce$\Gamma \fCenter !C$
  \RightLabel{\Intro{\lif}}
  \UnaryInf$\Gamma \fCenter !B \lif !C$
  \DisplayProof
  \]

  $\delta$'nın başka bir kuralla bittiği durumlar benzerdir.
\end{proof}

\begin{prop}
$\delta$, \Log{N2c} ya da \Log{N2i}'de $\Gamma \Sequent !A$'nın bir
!!a{proof}{}ı ise \Log{N1c}'de (\Log{N1i}'de) bir !!a{proof}~$\delta'$ vardır;
bunun son-!!{formula}{}ü $!A$'dır ve her $x:!B \in \Gamma$ için $!B^x$ açık
varsayımını taşır.
\end{prop}

\begin{proof}
  Yine $\delta$'nın !!{height}~değeri~$n$ üzerinde tümevarımla ilerliyoruz. $n=0$
  ise $\delta$, $x:!A \Sequent !A$ aksiyomudur. $\delta'$ !!{proof}{}ı $!A^x$
  varsayımıdır.

  $n>0$ ise $\delta$'daki son çıkarıma göre durumları ayırırız. Örneğin bu
  çıkarım~$\Intro{\lif}$ ise $\delta$ şöyle biter:
  \[
  \bottomAlignProof
  \AxiomC{}
  \RightLabel{$\delta_1$}
  \Deduce$x:!B, \Gamma \fCenter !C$
  \RightLabel{\Intro{\lif}}
  \UnaryInf$\Gamma \fCenter !B \lif !C$
  \DisplayProof
  \quad\text{ ya da }\quad
  \bottomAlignProof
  \AxiomC{}
  \RightLabel{$\delta_1$}
  \Deduce$\Gamma \fCenter !C$
  \RightLabel{\Intro{\lif}}
  \UnaryInf$\Gamma \fCenter !B \lif !C$
  \DisplayProof
  \]
  Tümevarım varsayımı gereği \Log{N1c}'de (\Log{N1i}'de) $!C$'nin bir
  $\delta_1'$ !!a{proof}{}ı vardır. Birinci durumda $!B^x$, $\delta_1'$'in açık
  bir varsayımıdır; ikinci durumda değildir. $\delta'$ !!{proof}{}ı olarak
  sırasıyla şunları alabiliriz:
  \[
  \bottomAlignProof
  \AxiomC{$\Discharge{!B}{x}$}
  \RightLabel{$\delta_1'$}
  \DeduceC{$!C$}
  \DischargeRule{\Intro{\lif}}{x}
  \UnaryInfC{$!B \lif !C$}
  \DisplayProof
  \quad\text{ ya da }\quad
  \bottomAlignProof
  \AxiomC{}
  \RightLabel{$\delta_1'$}
  \DeduceC{$!C$}
  \RightLabel{\Intro{\lif}}
  \UnaryInfC{$!B \lif !C$}
  \DisplayProof
  \]
\end{proof}

\end{document}
